\answer{OR16}
从任意未吸收状态，接下来的 $a+b$ 步若全向右便会吸收，其概率为 $2^{-(a+b)}$。因此 $\tau$ 有几何尾，所有阶矩均有限。

对称的独立增量使下列过程成为鞅：
\[
 S_n,\qquad S_n^2-n,\qquad S_n^3-3nS_n,
\]
\[
 S_n^4-6nS_n^2+3n^2+2n.
\]
例如分别展开 $(S_n\pm1)^j$ 再取平均即可核实。先在有界停时 $\tau\wedge N$ 应用可选停止，再令 $N\to\infty$。停止后的位置绝对值不超过 $\max(a,b)$，且 $\tau$ 的矩有限，故支配收敛允许这一极限。

设 $r=\Prob(S_\tau=b)$。前两个鞅分别给出
\[
 br-a(1-r)=0,\qquad
 \E\tau=\E S_\tau^2=b^2r+a^2(1-r),
\]
所以
\[
 \boxed{r=\frac a{a+b}},\qquad \boxed{\E\tau=ab}.
\]
为求二阶矩，令 $A=\E[\tau\ind_{\{S_\tau=b\}}]$、$B=\E[\tau\ind_{\{S_\tau=-a\}}]$。三次鞅给出
\[
 A+B=ab,\qquad
 bA-aB=\frac{\E S_\tau^3}{3}=\frac{ab(b-a)}3.
\]
解这两个线性方程，得
\[
 A=\frac{ab(2a+b)}{3(a+b)},\qquad
 B=\frac{ab(a+2b)}{3(a+b)}.
\]
于是
\[
 \E(\tau S_\tau^2)=b^2A+a^2B
       =\frac{ab(a^2+ab+b^2)}3,
\]
\[
 \E S_\tau^4=ab(a^2-ab+b^2).
\]
代入四次鞅的停止等式
\[
 0=\E S_\tau^4-6\E(\tau S_\tau^2)
       +3\E\tau^2+2\E\tau,
\]
最终得到
\[
 \boxed{\E\tau^2=\frac{ab(a^2+3ab+b^2-2)}3}.
\]
